\setcounter{chapter}{31}
\chapter{Routing}\label{ch:32-routing}

\chapterrule

\hyperref[ch:31-input-handling]{Chapter 31} validated the incoming request. Before validation can even begin, however, Flask must determine \emph{which handler function} the request belongs to. That determination is \textbf{routing}: the process of matching a URL and HTTP method to a Python function.

Flask's routing system is provided by \textbf{Werkzeug}, the library that Flask is built on top of. Werkzeug's routing is sophisticated~--- it handles URL patterns with variables, HTTP method constraints, custom type converters, and detailed error responses~--- but it is also fully transparent once you know where to look.

\begin{objectives}

By the end of this chapter, you will understand:

\begin{itemize}
\tightlist
\item
  How \texttt{@app.route()} registers a route at import time
\item
  What the Werkzeug \texttt{Map} and \texttt{Rule} objects are and what they contain
\item
  How URL matching works: the algorithm, the priority rules, the edge cases
\item
  What URL variable converters are and how to define custom ones
\item
  How 404 (Not Found) and 405 (Method Not Allowed) responses are generated
\item
  How Blueprint routes integrate into the application URL map
\item
  How to inspect the registered routes in a running application
\end{itemize}

\end{objectives}

\setcounter{section}{0}
\section{\texorpdfstring{Registration: What \texttt{@app.route()} Does}{Registration: What @app.route() Does}}\label{32-routing:321-registration-what-approute-does}

When you write:

\begin{codebox}[short]

\begin{Shaded}
\begin{Highlighting}[]
\AttributeTok{@app.route}\NormalTok{(}\StringTok{"/notes/"}\NormalTok{, methods}\OperatorTok{=}\NormalTok{[}\StringTok{"GET"}\NormalTok{])}
\KeywordTok{def}\NormalTok{ list\_notes():}
\NormalTok{    ...}
\end{Highlighting}
\end{Shaded}

\end{codebox}

\noindent{}this is Python decorator syntax. It is exactly equivalent to:

\begin{codebox}[short]

\begin{Shaded}
\begin{Highlighting}[]
\KeywordTok{def}\NormalTok{ list\_notes():}
\NormalTok{    ...}

\NormalTok{app.route(}\StringTok{"/notes/"}\NormalTok{, methods}\OperatorTok{=}\NormalTok{[}\StringTok{"GET"}\NormalTok{])(list\_notes)}
\end{Highlighting}
\end{Shaded}

\end{codebox}

\noindent{}Which is exactly equivalent to:

\begin{codebox}[short]

\begin{Shaded}
\begin{Highlighting}[]
\KeywordTok{def}\NormalTok{ list\_notes():}
\NormalTok{    ...}

\NormalTok{app.add\_url\_rule(}\StringTok{"/notes/"}\NormalTok{, endpoint}\OperatorTok{=}\StringTok{"list\_notes"}\NormalTok{, view\_func}\OperatorTok{=}\NormalTok{list\_notes, methods}\OperatorTok{=}\NormalTok{[}\StringTok{"GET"}\NormalTok{])}
\end{Highlighting}
\end{Shaded}

\end{codebox}

\noindent{}\texttt{add\_url\_rule} is the actual method that does the work. The decorator is syntactic sugar.

\subsection*{When does registration happen?}\phantomsection\label{32-routing:when-does-registration-happen}

Route registration happens at \textbf{import time}~--- when the module containing the \texttt{@app.route()} decorator is imported by Python. It does \emph{not} happen when a request arrives.

\begin{outputbox}
\begin{Verbatim}[fontsize=\codesize, breaklines, breakanywhere, breaksymbolleft={\tiny\textcolor{muted}{$\hookrightarrow$}}]
Timeline (production: gunicorn "app:create_app()"; in development,
python run.py does the same through run.py):

  Gunicorn worker imports the app package (app/__init__.py)
      ↓
  Gunicorn calls create_app()
      ↓
  create_app() imports and registers the notes Blueprint
      ↓
  Python imports app/routes/notes.py
      ↓
  @notes_bp.route("/") decorator executes → Rule added to Blueprint
  @notes_bp.route("/<int:note_id>") decorator executes → Rule added
      ↓
  create_app() calls app.register_blueprint(notes_bp)
      ↓
  Blueprint's rules are copied into app.url_map
      ↓
  URL map is complete. Ready to match requests.

  -----------------------------------------------

  First request arrives → the URL map is matched → handler found quickly
\end{Verbatim}
\end{outputbox}

\noindent{}The URL map is built once at startup and then used for every subsequent request. This is intentional~--- routing lookup must be fast, and building the map only once makes it fast.

\subsection*{\texorpdfstring{What \texttt{add\_url\_rule} does internally}{What add\_url\_rule does internally}}\phantomsection\label{32-routing:what-add-url-rule-does-internally}

\begin{codebox}

\begin{Shaded}
\begin{Highlighting}[]
\CommentTok{\# Flask\textquotesingle{}s add\_url\_rule (simplified to show the concepts)}
\KeywordTok{def}\NormalTok{ add\_url\_rule(}\VariableTok{self}\NormalTok{, rule, endpoint}\OperatorTok{=}\VariableTok{None}\NormalTok{, view\_func}\OperatorTok{=}\VariableTok{None}\NormalTok{, methods}\OperatorTok{=}\VariableTok{None}\NormalTok{, }\OperatorTok{**}\NormalTok{options):}
    \CommentTok{\# 1. Determine the endpoint name}
    \ControlFlowTok{if}\NormalTok{ endpoint }\KeywordTok{is} \VariableTok{None}\NormalTok{:}
\NormalTok{        endpoint }\OperatorTok{=}\NormalTok{ view\_func.}\VariableTok{\_\_name\_\_}  \CommentTok{\# "list\_notes"}

    \CommentTok{\# 2. Determine which HTTP methods this rule handles}
    \ControlFlowTok{if}\NormalTok{ methods }\KeywordTok{is} \VariableTok{None}\NormalTok{:}
\NormalTok{        methods }\OperatorTok{=} \BuiltInTok{getattr}\NormalTok{(view\_func, }\StringTok{"methods"}\NormalTok{, }\VariableTok{None}\NormalTok{) }\KeywordTok{or}\NormalTok{ (}\StringTok{"GET"}\NormalTok{,)}
\NormalTok{    methods }\OperatorTok{=}\NormalTok{ \{m.upper() }\ControlFlowTok{for}\NormalTok{ m }\KeywordTok{in}\NormalTok{ methods\}  }\CommentTok{\# Normalize to uppercase}

    \CommentTok{\# 3. Add HEAD wherever GET is allowed (Werkzeug\textquotesingle{}s Rule does this:}
    \CommentTok{\#    a HEAD request is a GET without the body)}
    \ControlFlowTok{if} \StringTok{"GET"} \KeywordTok{in}\NormalTok{ methods:}
\NormalTok{        methods.add(}\StringTok{"HEAD"}\NormalTok{)}

    \CommentTok{\# 4. Add OPTIONS (Flask answers it automatically, unless the view}
    \CommentTok{\#    sets provide\_automatic\_options = False)}
\NormalTok{    methods.add(}\StringTok{"OPTIONS"}\NormalTok{)}

    \CommentTok{\# 5. Create a Werkzeug Rule object}
\NormalTok{    rule\_obj }\OperatorTok{=} \VariableTok{self}\NormalTok{.url\_rule\_class(}
\NormalTok{        rule,            }\CommentTok{\# The URL pattern: "/notes/"}
\NormalTok{        methods}\OperatorTok{=}\NormalTok{methods, }\CommentTok{\# \{"GET", "HEAD", "OPTIONS"\}}
\NormalTok{        endpoint}\OperatorTok{=}\NormalTok{endpoint,}
\NormalTok{    )}

    \CommentTok{\# 6. Add the Rule to the URL Map}
    \VariableTok{self}\NormalTok{.url\_map.add(rule\_obj)}

    \CommentTok{\# 7. Register the view function for this endpoint}
    \VariableTok{self}\NormalTok{.view\_functions[endpoint] }\OperatorTok{=}\NormalTok{ view\_func}
\end{Highlighting}
\end{Shaded}

\end{codebox}

\noindent{}After registration, Flask has:

\begin{enumerate}
\def\labelenumi{\arabic{enumi}.}
\tightlist
\item
  A \texttt{url\_map} containing \texttt{Rule} objects (URL patterns)
\item
  A \texttt{view\_functions} dict mapping endpoint names to callable functions
\end{enumerate}

\noindent{}When a request arrives, Flask uses \texttt{url\_map} to find the endpoint name, then uses \texttt{view\_functions} to find the callable.

\setcounter{section}{1}
\section{The URL Map and Rule Objects}\label{32-routing:322-the-url-map-and-rule-objects}

\subsection*{\texorpdfstring{Werkzeug's \texttt{Map}}{Werkzeug's Map}}\phantomsection\label{32-routing:werkzeugs-map}

The \texttt{Map} is a collection of \texttt{Rule} objects. You can inspect it directly:

\begin{codebox}[short]

\begin{Shaded}
\begin{Highlighting}[]
\CommentTok{\# Run this in a Flask shell: flask shell}
\ImportTok{from}\NormalTok{ app }\ImportTok{import}\NormalTok{ create\_app}
\NormalTok{app }\OperatorTok{=}\NormalTok{ create\_app()}

\CommentTok{\# Print all registered rules}
\ControlFlowTok{for}\NormalTok{ rule }\KeywordTok{in}\NormalTok{ app.url\_map.iter\_rules():}
    \BuiltInTok{print}\NormalTok{(}\SpecialStringTok{f"}\SpecialCharTok{\{}\NormalTok{rule}\SpecialCharTok{.}\NormalTok{endpoint}\SpecialCharTok{:30s\}}\SpecialStringTok{  }\SpecialCharTok{\{}\StringTok{\textquotesingle{},\textquotesingle{}}\SpecialCharTok{.}\NormalTok{join(}\BuiltInTok{sorted}\NormalTok{(rule.methods))}\SpecialCharTok{:30s\}}\SpecialStringTok{  }\SpecialCharTok{\{}\NormalTok{rule}\SpecialCharTok{.}\NormalTok{rule}\SpecialCharTok{\}}\SpecialStringTok{"}\NormalTok{)}
\end{Highlighting}
\end{Shaded}

\end{codebox}

\noindent{}For the Notes API as it stands after \hyperref[ch:16-logging-error-handling]{Chapter 16}, this produces (with the column widths narrowed to fit the page):

\begin{outputbox}[short]
\begin{Verbatim}[fontsize=\codesize, breaklines, breakanywhere, breaksymbolleft={\tiny\textcolor{muted}{$\hookrightarrow$}}]
static                  GET,HEAD,OPTIONS        /static/<path:filename>
health.health           GET,HEAD,OPTIONS        /health
notes.list_notes        GET,HEAD,OPTIONS        /notes/
notes.create_note       OPTIONS,POST            /notes/
notes.get_note          GET,HEAD,OPTIONS        /notes/<int:note_id>
notes.delete_note       DELETE,OPTIONS          /notes/<int:note_id>
\end{Verbatim}
\end{outputbox}

\noindent{}Several things are apparent:

\begin{itemize}
\tightlist
\item
  Each route has an \textbf{endpoint} name (e.g., \texttt{notes.list\_notes}~--- the Blueprint name prefixed)
\item
  Methods are automatically extended with \texttt{OPTIONS}, and with \texttt{HEAD} wherever \texttt{GET} is allowed
\item
  Flask registers a \texttt{static} route automatically for serving static files
\item
  Multiple rules can share the same URL pattern but differ by HTTP method
\end{itemize}

\subsection*{\texorpdfstring{Werkzeug's \texttt{Rule}}{Werkzeug's Rule}}\phantomsection\label{32-routing:werkzeugs-rule}

A \texttt{Rule} is the unit of URL matching. Each \texttt{Rule} contains:

\begin{codebox}[short]

\begin{Shaded}
\begin{Highlighting}[]
\NormalTok{rule.rule      }\CommentTok{\# The URL pattern string: "/notes/\textless{}int:note\_id\textgreater{}"}
\NormalTok{rule.methods   }\CommentTok{\# The set of HTTP methods: \{"GET", "HEAD", "OPTIONS"\}}
\NormalTok{rule.endpoint  }\CommentTok{\# The endpoint name: "notes.get\_note"}
\end{Highlighting}
\end{Shaded}

\end{codebox}

\noindent{}Conceptually, each \texttt{Rule} is a pattern that a path either matches or not, and the \texttt{\textless{}int:note\_id\textgreater{}} part is a \textbf{converter}~--- it specifies both the pattern for that segment (the regular expression \texttt{\textbackslash{}d+}) and the Python type to convert the matched string to (\texttt{int}):

\begin{outputbox}[short]
\begin{Verbatim}[fontsize=\codesize, breaklines, breakanywhere, breaksymbolleft={\tiny\textcolor{muted}{$\hookrightarrow$}}]
"/notes/<int:note_id>"
         ↓ conceptually
^/notes/(?P<note_id>\d+)$
\end{Verbatim}
\end{outputbox}

\noindent{}(Older Werkzeug versions compiled each rule to one regular expression like this and tried the rules in turn. Since Werkzeug 2.2, the \texttt{Map} compiles all rules together into a \textbf{state machine} that walks the path one segment at a time~--- \texttt{notes}, then \texttt{42}~--- so matching stays fast however many routes there are. The outcome is the same; the trace below shows the conceptual version.)

\setcounter{section}{2}
\section{URL Matching: How a Request Finds Its Handler}\label{32-routing:323-url-matching-how-a-request-finds-its-handler}

When a request arrives, Flask calls:

\begin{codebox}[short]

\begin{Shaded}
\begin{Highlighting}[]
\NormalTok{adapter }\OperatorTok{=}\NormalTok{ app.url\_map.bind\_to\_environ(environ)}
\NormalTok{endpoint, values }\OperatorTok{=}\NormalTok{ adapter.match()}
\end{Highlighting}
\end{Shaded}

\end{codebox}

\noindent{}Let's trace this for \texttt{GET\ /notes/42}:

\begin{outputbox}
\begin{Verbatim}[fontsize=\codesize, breaklines, breakanywhere, breaksymbolleft={\tiny\textcolor{muted}{$\hookrightarrow$}}]
Step 1: bind_to_environ
   Creates a MapAdapter bound to the current request's host and scheme.
   The adapter knows: host="api.notes.example.com", scheme="https"

Step 2: adapter.match()  (conceptually)
   Checking Rule: "/notes/"
       "/notes/42" has a second segment "42" → no match

   Checking Rule: "/notes/<int:note_id>"  (segment pattern: \d+)
       "/notes/42" → MATCH
       Captured: {"note_id": "42"}

Step 3: Method check
   The matched Rule handles: {"GET", "HEAD", "OPTIONS"}
   The request method is: "GET"
   → GET is in the allowed set → proceed

Step 4: Type conversion
   "42" is passed to the int converter → 42 (Python int)

Step 5: Return
   endpoint = "notes.get_note"
   values   = {"note_id": 42}
\end{Verbatim}
\end{outputbox}

\noindent{}Flask then looks up \texttt{view\_}\allowbreak{}\texttt{functions{[}"notes.}\allowbreak{}\texttt{get\_}\allowbreak{}\texttt{note"{]}} and calls it with \texttt{note\_id=42}.

\subsection*{What happens when there is no match}\phantomsection\label{32-routing:what-happens-when-there-is-no-match}

\textbf{Path not found}~--- no Rule matches the URL path at all:

\begin{codebox}[short]

\begin{Shaded}
\begin{Highlighting}[]
\CommentTok{\# GET /nonexistent/}

\CommentTok{\# adapter.match() raises: werkzeug.exceptions.NotFound}
\CommentTok{\# Flask catches this and calls: app.handle\_http\_exception(NotFound())}
\CommentTok{\# Which calls the registered 404 error handler}
\CommentTok{\# Result: HTTP 404 response}
\end{Highlighting}
\end{Shaded}

\end{codebox}

\noindent{}\textbf{Method not allowed}~--- a Rule matches the path, but not for this HTTP method:

\begin{codebox}[short]

\begin{Shaded}
\begin{Highlighting}[]
\CommentTok{\# DELETE /notes/  (there is a GET rule for /notes/ but not DELETE)}

\CommentTok{\# adapter.match() raises:}
\CommentTok{\#   werkzeug.exceptions.MethodNotAllowed(valid\_methods=["GET", "HEAD", "OPTIONS", "POST"])}
\CommentTok{\# (both rules for /notes/ count: list\_notes and create\_note)}
\CommentTok{\# Flask calls the 405 error handler, which returns HTTP 405 with the header}
\CommentTok{\#   Allow: GET, HEAD, OPTIONS, POST}
\end{Highlighting}
\end{Shaded}

\end{codebox}

\noindent{}The distinction between 404 and 405 matters:

\begin{itemize}
\tightlist
\item
  \textbf{404}: The URL doesn't exist. The client should not retry with a different method.
\item
  \textbf{405}: The URL exists but the method is wrong. The \texttt{Allow} header tells the client what methods work.
\end{itemize}

\noindent{}This distinction is why methods belong in the route declaration (\texttt{methods={[}...{]}}) rather than in \texttt{if\ request.}\allowbreak{}\texttt{method\ =}\allowbreak{}\texttt{=}\allowbreak{}\texttt{\ .}\allowbreak{}\texttt{.}\allowbreak{}\texttt{.} branches inside a view: only declared methods let Flask answer 405 with a correct \texttt{Allow} header.

\setcounter{section}{3}
\section{URL Variable Converters}\label{32-routing:324-url-variable-converters}

Converters are the \texttt{\textless{}type:name\textgreater{}} syntax in URL patterns. They perform two jobs:

\begin{enumerate}
\def\labelenumi{\arabic{enumi}.}
\tightlist
\item
  \textbf{Pattern}: control what part of the URL matches (via regex)
\item
  \textbf{Conversion}: transform the matched string into a Python value
\end{enumerate}

\Needspace{19\baselineskip}

\subsection*{Built-in converters}\phantomsection\label{32-routing:built-in-converters}

{\def\LTcaptype{none} % do not increment counter
\begin{longtable}[]{@{}
  >{\raggedright\arraybackslash}p{(\linewidth - 6\tabcolsep) * \real{0.1481}}
  >{\raggedright\arraybackslash}p{(\linewidth - 6\tabcolsep) * \real{0.2251}}
  >{\raggedright\arraybackslash}p{(\linewidth - 6\tabcolsep) * \real{0.1709}}
  >{\raggedright\arraybackslash}p{(\linewidth - 6\tabcolsep) * \real{0.4558}}@{}}
\toprule\noalign{}
\begin{minipage}[b]{\linewidth}\raggedright
Converter
\end{minipage} & \begin{minipage}[b]{\linewidth}\raggedright
Pattern
\end{minipage} & \begin{minipage}[b]{\linewidth}\raggedright
Python type
\end{minipage} & \begin{minipage}[b]{\linewidth}\raggedright
Example
\end{minipage} \\
\midrule\noalign{}
\endhead
\bottomrule\noalign{}
\endlastfoot
\texttt{string} & \texttt{{[}\^{}/{]}+} & \texttt{str} & \texttt{\textless{}string:name\textgreater{}} matches \texttt{foo} \\
\texttt{int} & \texttt{{[}0-9{]}+} & \texttt{int} & \texttt{\textless{}int:id\textgreater{}} matches \texttt{42} \\
\texttt{float} & \texttt{{[}0-9{]}+\textbackslash{}.{[}0-9{]}+} & \texttt{float} & \texttt{\textless{}float:ratio\textgreater{}} matches \texttt{3.14} \\
\texttt{path} & \texttt{.+} & \texttt{str} & \texttt{\textless{}path:filename\textgreater{}} matches \texttt{dir/file.txt} \\
\texttt{uuid} & UUID pattern & \texttt{uuid.UUID} & \texttt{\textless{}uuid:id\textgreater{}} matches \texttt{550e8400-...} \\
\end{longtable}
}

The default converter (used when you write just \texttt{\textless{}name\textgreater{}}) is \texttt{string}.

\begin{codebox}[short]

\begin{Shaded}
\begin{Highlighting}[]
\CommentTok{\# These are equivalent:}
\AttributeTok{@app.route}\NormalTok{(}\StringTok{"/users/\textless{}name\textgreater{}"}\NormalTok{)}
\AttributeTok{@app.route}\NormalTok{(}\StringTok{"/users/\textless{}string:name\textgreater{}"}\NormalTok{)}

\CommentTok{\# int converter: only matches digits, converts to int}
\AttributeTok{@app.route}\NormalTok{(}\StringTok{"/notes/\textless{}int:note\_id\textgreater{}"}\NormalTok{)}
\CommentTok{\# /notes/42  → note\_id=42 (int)}
\CommentTok{\# /notes/abc → 404 Not Found (no match)}

\CommentTok{\# path converter: matches slashes}
\AttributeTok{@app.route}\NormalTok{(}\StringTok{"/files/\textless{}path:filename\textgreater{}"}\NormalTok{)}
\CommentTok{\# /files/images/photo.jpg → filename="images/photo.jpg"}
\end{Highlighting}
\end{Shaded}

\end{codebox}

\subsection*{\texorpdfstring{Why \texttt{\textless{}int:note\_id\textgreater{}} is better than
\texttt{\textless{}note\_id\textgreater{}}}{Why \textless int:note\_id\textgreater{} is better than \textless note\_id\textgreater{}}}\phantomsection\label{32-routing:why-intnote-id-is-better-than-note-id}

Consider what happens with each approach when someone requests \texttt{/notes/abc}:

\begin{codebox}[short]

\begin{Shaded}
\begin{Highlighting}[]
\CommentTok{\# With string converter (default):}
\AttributeTok{@app.route}\NormalTok{(}\StringTok{"/notes/\textless{}note\_id\textgreater{}"}\NormalTok{)}
\KeywordTok{def}\NormalTok{ get\_note(note\_id):}
    \CommentTok{\# note\_id = "abc" (a string)}
    \CommentTok{\# You must validate and convert it yourself}
    \ControlFlowTok{try}\NormalTok{:}
\NormalTok{        note\_id }\OperatorTok{=} \BuiltInTok{int}\NormalTok{(note\_id)}
    \ControlFlowTok{except} \PreprocessorTok{ValueError}\NormalTok{:}
        \ControlFlowTok{return}\NormalTok{ \{}\StringTok{"error"}\NormalTok{: }\StringTok{"note\_id must be an integer"}\NormalTok{\}, }\DecValTok{400}
\NormalTok{    ...}

\CommentTok{\# With int converter:}
\AttributeTok{@app.route}\NormalTok{(}\StringTok{"/notes/\textless{}int:note\_id\textgreater{}"}\NormalTok{)}
\KeywordTok{def}\NormalTok{ get\_note(note\_id):}
    \CommentTok{\# note\_id = 42 (already an int, guaranteed)}
    \CommentTok{\# /notes/abc never reaches here — 404 before the function is called}
\NormalTok{    ...}
\end{Highlighting}
\end{Shaded}

\end{codebox}

\noindent{}The \texttt{int} converter means \texttt{/notes/abc} returns 404 automatically~--- ``no such route''~--- which is correct. There is no note at \texttt{/notes/abc}; that URL doesn't exist.

\subsection*{Defining a custom converter}\phantomsection\label{32-routing:defining-a-custom-converter}

If the built-in converters don't cover your case, you can write your own:

\begin{codebox}

\begin{Shaded}
\begin{Highlighting}[]
\CommentTok{\# app/converters.py}
\ImportTok{from}\NormalTok{ werkzeug.routing }\ImportTok{import}\NormalTok{ BaseConverter}

\KeywordTok{class}\NormalTok{ SlugConverter(BaseConverter):}
    \CommentTok{"""}
\CommentTok{    Matches URL{-}safe slugs: lowercase letters, digits, hyphens.}
\CommentTok{    Example: /articles/my{-}first{-}post}
\CommentTok{    """}
    \CommentTok{\# This regex is what the converter matches in the URL}
    \CommentTok{\# (a single character, or several that do not start or end with "{-}")}
\NormalTok{    regex }\OperatorTok{=} \VerbatimStringTok{r"}\PreprocessorTok{[a{-}z0{-}9]}\NormalTok{(?:}\PreprocessorTok{[a{-}z0{-}9}\CharTok{\textbackslash{}{-}}\PreprocessorTok{]}\OperatorTok{*}\PreprocessorTok{[a{-}z0{-}9]}\NormalTok{)}\OperatorTok{?}\VerbatimStringTok{"}

    \KeywordTok{def}\NormalTok{ to\_python(}\VariableTok{self}\NormalTok{, value: }\BuiltInTok{str}\NormalTok{) }\OperatorTok{{-}\textgreater{}} \BuiltInTok{str}\NormalTok{:}
        \CommentTok{"""Called when matching: convert the URL string to a Python value."""}
        \ControlFlowTok{return}\NormalTok{ value  }\CommentTok{\# Already the right type}

    \KeywordTok{def}\NormalTok{ to\_url(}\VariableTok{self}\NormalTok{, value: }\BuiltInTok{str}\NormalTok{) }\OperatorTok{{-}\textgreater{}} \BuiltInTok{str}\NormalTok{:}
        \CommentTok{"""Called when building URLs: convert a Python value to a URL string."""}
        \ControlFlowTok{return}\NormalTok{ value  }\CommentTok{\# Build URLs from slugs exactly as they are}
\end{Highlighting}
\end{Shaded}

\end{codebox}

\noindent{}Register the converter in \texttt{create\_app()}:

\begin{codebox}[short]

\begin{Shaded}
\begin{Highlighting}[]
\CommentTok{\# app/\_\_init\_\_.py}
\ImportTok{from}\NormalTok{ app.converters }\ImportTok{import}\NormalTok{ SlugConverter}

\KeywordTok{def}\NormalTok{ create\_app():}
\NormalTok{    app }\OperatorTok{=}\NormalTok{ Flask(}\VariableTok{\_\_name\_\_}\NormalTok{)}
    \CommentTok{\# Register the converter BEFORE any blueprint that uses \textless{}slug:...\textgreater{}:}
    \CommentTok{\# the rules are built when the blueprint is registered, and an}
    \CommentTok{\# unknown converter name is an error at that moment}
\NormalTok{    app.url\_map.converters[}\StringTok{"slug"}\NormalTok{] }\OperatorTok{=}\NormalTok{ SlugConverter}
    \CommentTok{\# ... rest of create\_app (register\_blueprint calls come after)}
    \ControlFlowTok{return}\NormalTok{ app}
\end{Highlighting}
\end{Shaded}

\end{codebox}

\noindent{}Use it in routes:

\begin{codebox}[short]

\begin{Shaded}
\begin{Highlighting}[]
\AttributeTok{@bp.route}\NormalTok{(}\StringTok{"/articles/\textless{}slug:article\_slug\textgreater{}"}\NormalTok{)}
\KeywordTok{def}\NormalTok{ get\_article(article\_slug: }\BuiltInTok{str}\NormalTok{):}
    \CommentTok{\# /articles/my{-}first{-}post → article\_slug="my{-}first{-}post"}
    \CommentTok{\# /articles/CAPS → 404 (doesn\textquotesingle{}t match regex)}
\NormalTok{    ...}
\end{Highlighting}
\end{Shaded}

\end{codebox}

\setcounter{section}{4}
\section{Routing Priority and Ambiguity}\label{32-routing:325-routing-priority-and-ambiguity}

When multiple Rules could match the same URL, Werkzeug uses \textbf{specificity ordering} to pick the winner.

\subsection*{Specificity rules}\phantomsection\label{32-routing:specificity-rules}

More specific rules take priority over less specific ones:

\begin{codebox}[short]

\begin{Shaded}
\begin{Highlighting}[]
\CommentTok{\# Rule 1: Specific static segment}
\AttributeTok{@app.route}\NormalTok{(}\StringTok{"/notes/new"}\NormalTok{)}

\CommentTok{\# Rule 2: Variable segment that accepts any text}
\AttributeTok{@app.route}\NormalTok{(}\StringTok{"/notes/\textless{}name\textgreater{}"}\NormalTok{)          }\CommentTok{\# \textless{}name\textgreater{} = \textless{}string:name\textgreater{}}

\CommentTok{\# GET /notes/new:}
\CommentTok{\#   Rule 1 matches: endpoint="new\_note\_form"}
\CommentTok{\#   Rule 2 ALSO matches (name="new") — both rules fit}
\CommentTok{\#   Rule 1 wins because static segments are more specific than variable segments}

\CommentTok{\# GET /notes/groceries:}
\CommentTok{\#   Rule 1 doesn\textquotesingle{}t match ("groceries" ≠ "new")}
\CommentTok{\#   Rule 2 matches: name="groceries"}
\end{Highlighting}
\end{Shaded}

\end{codebox}

\noindent{}Static path segments always beat variable segments. \texttt{/notes/new} is more specific than \texttt{/}\allowbreak{}\texttt{notes/}\allowbreak{}\texttt{\textless{}int:note\_id\textgreater{}}.

\subsection*{Trailing slash handling}\phantomsection\label{32-routing:trailing-slash-handling}

Flask and Werkzeug have a specific rule about trailing slashes that catches many developers off guard.

\textbf{Rule with trailing slash}: \texttt{@app.route("/}\allowbreak{}\texttt{notes/}\allowbreak{}\texttt{")}

\begin{itemize}
\tightlist
\item
  GET \texttt{/notes/} → matches (200 OK)
\item
  GET \texttt{/notes} → redirect to \texttt{/notes/} (\textbf{308} Permanent Redirect; older Werkzeug versions used 301)
\end{itemize}

\noindent{}\textbf{Rule without trailing slash}: \texttt{@app.route("/}\allowbreak{}\texttt{notes")}

\begin{itemize}
\tightlist
\item
  GET \texttt{/notes} → matches (200 OK)
\item
  GET \texttt{/notes/} → 404 Not Found
\end{itemize}

\noindent{}The behavior with a trailing slash mimics how web servers handle directory URLs. The 308 status matters for an API: unlike 301, it tells the client to repeat the \emph{same} method with the \emph{same} body, so a \texttt{POST\ /notes} is re-sent as \texttt{POST\ /notes/}. But the client must follow redirects at all~--- \texttt{curl} does only with \texttt{-L}, and some HTTP libraries do not re-send request bodies~--- so for API clients a redirect is a trap, not a convenience. The Notes API uses a trailing slash for its collection endpoint (\texttt{/notes/}) and none for item endpoints (\texttt{/notes/42}), the convention Flask's documentation suggests, and documents the exact URLs so clients never depend on the redirect. (Many JSON-API style guides choose the opposite, no trailing slashes anywhere; either works if it is consistent.)

\begin{codebox}[short]

\begin{Shaded}
\begin{Highlighting}[]
\CommentTok{\# Collection endpoints: trailing slash (so /notes redirects to /notes/)}
\AttributeTok{@notes\_bp.route}\NormalTok{(}\StringTok{"/"}\NormalTok{, methods}\OperatorTok{=}\NormalTok{[}\StringTok{"GET"}\NormalTok{])        }\CommentTok{\# GET /notes/}
\AttributeTok{@notes\_bp.route}\NormalTok{(}\StringTok{"/"}\NormalTok{, methods}\OperatorTok{=}\NormalTok{[}\StringTok{"POST"}\NormalTok{])       }\CommentTok{\# POST /notes/}

\CommentTok{\# Item endpoints: no trailing slash}
\AttributeTok{@notes\_bp.route}\NormalTok{(}\StringTok{"/\textless{}int:note\_id\textgreater{}"}\NormalTok{, methods}\OperatorTok{=}\NormalTok{[}\StringTok{"GET"}\NormalTok{])     }\CommentTok{\# GET /notes/42}
\AttributeTok{@notes\_bp.route}\NormalTok{(}\StringTok{"/\textless{}int:note\_id\textgreater{}"}\NormalTok{, methods}\OperatorTok{=}\NormalTok{[}\StringTok{"PUT"}\NormalTok{])     }\CommentTok{\# PUT /notes/42 (Chapter 33)}
\AttributeTok{@notes\_bp.route}\NormalTok{(}\StringTok{"/\textless{}int:note\_id\textgreater{}"}\NormalTok{, methods}\OperatorTok{=}\NormalTok{[}\StringTok{"DELETE"}\NormalTok{])  }\CommentTok{\# DELETE /notes/42}
\end{Highlighting}
\end{Shaded}

\end{codebox}

\setcounter{section}{5}
\section{Blueprints and the URL Map}\label{32-routing:326-blueprints-and-the-url-map}

A \textbf{Blueprint} is a collection of routes that can be registered on an app. Blueprints have their own route registry, separate from the app's \texttt{url\_map}. When you call \texttt{app.}\allowbreak{}\texttt{register\_}\allowbreak{}\texttt{blueprint(bp)}, the Blueprint's rules are copied into the app's \texttt{url\_map}.

\subsection*{The url\_prefix}\phantomsection\label{32-routing:the-url-prefix}

\begin{codebox}[short]

\begin{Shaded}
\begin{Highlighting}[]
\CommentTok{\# In app/routes/notes.py:}
\NormalTok{notes\_bp }\OperatorTok{=}\NormalTok{ Blueprint(}\StringTok{"notes"}\NormalTok{, }\VariableTok{\_\_name\_\_}\NormalTok{, url\_prefix}\OperatorTok{=}\StringTok{"/notes"}\NormalTok{)}

\CommentTok{\# In app/\_\_init\_\_.py:}
\NormalTok{app.register\_blueprint(notes\_bp)}

\AttributeTok{@notes\_bp.route}\NormalTok{(}\StringTok{"/"}\NormalTok{)              }\CommentTok{\# The Blueprint{-}relative path}
\KeywordTok{def}\NormalTok{ list\_notes():}
\NormalTok{    ...}

\AttributeTok{@notes\_bp.route}\NormalTok{(}\StringTok{"/\textless{}int:note\_id\textgreater{}"}\NormalTok{)  }\CommentTok{\# Blueprint{-}relative}
\KeywordTok{def}\NormalTok{ get\_note(note\_id):}
\NormalTok{    ...}
\end{Highlighting}
\end{Shaded}

\end{codebox}

\noindent{}After registration, the actual URL patterns in the app's \texttt{url\_map} are:

\begin{itemize}
\tightlist
\item
  \texttt{/notes} + \texttt{/} = \texttt{/notes/}
\item
  \texttt{/notes} + \texttt{/\textless{}int:note\_id\textgreater{}} = \texttt{/}\allowbreak{}\texttt{notes/}\allowbreak{}\texttt{\textless{}int:note\_id\textgreater{}}
\end{itemize}

\noindent{}The \texttt{url\_prefix} acts as a namespace for all routes in the Blueprint. (It can also be passed to \texttt{register\_}\allowbreak{}\texttt{blueprint(notes\_}\allowbreak{}\texttt{bp,}\allowbreak{}\texttt{\ url\_}\allowbreak{}\texttt{prefix=}\allowbreak{}\texttt{"/}\allowbreak{}\texttt{notes")} instead, which overrides the Blueprint's own; the Notes API sets it on the Blueprint, next to its routes.)

\subsection*{Blueprint endpoint names}\phantomsection\label{32-routing:blueprint-endpoint-names}

Blueprint route endpoint names are automatically prefixed with the Blueprint name:

\begin{codebox}[short]

\begin{Shaded}
\begin{Highlighting}[]
\CommentTok{\# Blueprint name: "notes"}
\CommentTok{\# Route function: list\_notes}

\CommentTok{\# Resulting endpoint name: "notes.list\_notes"}
\end{Highlighting}
\end{Shaded}

\end{codebox}

\noindent{}This matters when building URLs with \texttt{url\_for()}:

\begin{codebox}[short]

\begin{Shaded}
\begin{Highlighting}[]
\ImportTok{from}\NormalTok{ flask }\ImportTok{import}\NormalTok{ url\_for}

\CommentTok{\# Inside a request context:}
\NormalTok{url\_for(}\StringTok{"notes.list\_notes"}\NormalTok{)           }\CommentTok{\# "/notes/"}
\NormalTok{url\_for(}\StringTok{"notes.get\_note"}\NormalTok{, note\_id}\OperatorTok{=}\DecValTok{42}\NormalTok{) }\CommentTok{\# "/notes/42"}

\CommentTok{\# Without the Blueprint prefix, this raises BuildError:}
\NormalTok{url\_for(}\StringTok{"list\_notes"}\NormalTok{)  }\CommentTok{\# ERROR: no endpoint named "list\_notes"}
\end{Highlighting}
\end{Shaded}

\end{codebox}

\noindent{}The Blueprint prefix prevents endpoint name collisions between Blueprints. If two Blueprints both have a \texttt{list} function, their endpoints are \texttt{users.list} and \texttt{notes.list}~--- no conflict.

\setcounter{section}{6}
\section{\texorpdfstring{URL Building with \texttt{url\_for()}}{URL Building with url\_for()}}\label{32-routing:327-url-building-with-url-for}

\texttt{url\_for()} is the reverse of URL matching: given an endpoint name (and any required variables), it produces the URL string.

\begin{codebox}[short]

\begin{Shaded}
\begin{Highlighting}[]
\ImportTok{from}\NormalTok{ flask }\ImportTok{import}\NormalTok{ url\_for}

\CommentTok{\# Basic usage}
\NormalTok{url\_for(}\StringTok{"notes.list\_notes"}\NormalTok{)              }\CommentTok{\# "/notes/"}
\NormalTok{url\_for(}\StringTok{"notes.get\_note"}\NormalTok{, note\_id}\OperatorTok{=}\DecValTok{42}\NormalTok{)    }\CommentTok{\# "/notes/42"}
\NormalTok{url\_for(}\StringTok{"notes.delete\_note"}\NormalTok{, note\_id}\OperatorTok{=}\DecValTok{7}\NormalTok{)  }\CommentTok{\# "/notes/7"}

\CommentTok{\# Extra keyword arguments become query string parameters}
\NormalTok{url\_for(}\StringTok{"notes.list\_notes"}\NormalTok{, page}\OperatorTok{=}\DecValTok{2}\NormalTok{, limit}\OperatorTok{=}\DecValTok{10}\NormalTok{)  }\CommentTok{\# "/notes/?page=2\&limit=10"}

\CommentTok{\# \_external=True builds an absolute URL}
\NormalTok{url\_for(}\StringTok{"notes.get\_note"}\NormalTok{, note\_id}\OperatorTok{=}\DecValTok{42}\NormalTok{, \_external}\OperatorTok{=}\VariableTok{True}\NormalTok{)}
\CommentTok{\# "http://localhost:5000/notes/42" (in dev)}
\CommentTok{\# "https://api.notes.example.com/notes/42" (in production — the host comes from}
\CommentTok{\#   the request\textquotesingle{}s Host header, and the scheme from X{-}Forwarded{-}Proto via ProxyFix)}
\end{Highlighting}
\end{Shaded}

\end{codebox}

\subsection*{\texorpdfstring{Why use \texttt{url\_for()} instead of hardcoding URLs?}{Why use url\_for() instead of hardcoding URLs?}}\phantomsection\label{32-routing:why-use-url-for-instead-of-hardcoding-urls}

\begin{codebox}[short]

\begin{Shaded}
\begin{Highlighting}[]
\CommentTok{\# Fragile: hardcoded URL}
\ControlFlowTok{return}\NormalTok{ redirect(}\StringTok{"/notes/"} \OperatorTok{+} \BuiltInTok{str}\NormalTok{(note\_id))}

\CommentTok{\# Correct: use url\_for}
\ControlFlowTok{return}\NormalTok{ redirect(url\_for(}\StringTok{"notes.get\_note"}\NormalTok{, note\_id}\OperatorTok{=}\NormalTok{note\_id))}
\end{Highlighting}
\end{Shaded}

\end{codebox}

\noindent{}If you change the URL pattern in \texttt{@notes\_}\allowbreak{}\texttt{bp.}\allowbreak{}\texttt{route("/}\allowbreak{}\texttt{\textless{}int:}\allowbreak{}\texttt{note\_}\allowbreak{}\texttt{id\textgreater{}")} to something else, every hardcoded URL breaks. \texttt{url\_for()} recomputes the URL from the current rule, so it always stays in sync.

\subsection*{\texorpdfstring{\texttt{url\_for()} requires a request context}{url\_for() requires a request context}}\phantomsection\label{32-routing:url-for-requires-a-request-context}

\texttt{url\_for()} normally requires an active request context (or an application context with \texttt{SERVER\_NAME} configured~--- a setting best left unset behind a proxy, because it also makes Flask reject requests for any other host name). In test code that is not itself handling a request, push one:

\begin{codebox}[short]

\begin{Shaded}
\begin{Highlighting}[]
\ControlFlowTok{with}\NormalTok{ app.test\_request\_context():}
\NormalTok{    url }\OperatorTok{=}\NormalTok{ url\_for(}\StringTok{"notes.get\_note"}\NormalTok{, note\_id}\OperatorTok{=}\DecValTok{42}\NormalTok{)}
\end{Highlighting}
\end{Shaded}

\end{codebox}

\noindent{}(During a request made with the test client, the request context is pushed automatically~--- but only for the code that runs inside that request.)

\setcounter{section}{7}
\section{Inspecting Routes}\label{32-routing:328-inspecting-routes}

Several tools let you see all registered routes at runtime.

\subsection*{Flask CLI command}\phantomsection\label{32-routing:flask-cli-command}

\begin{codebox}[short]

\begin{Shaded}
\begin{Highlighting}[]
\ExtensionTok{flask} \AttributeTok{{-}{-}app}\NormalTok{ run routes}

\CommentTok{\# Output:}
\ExtensionTok{Endpoint}\NormalTok{           Methods  Rule}
\ExtensionTok{{-}{-}{-}{-}{-}{-}{-}{-}{-}{-}{-}{-}{-}{-}{-}{-}{-}}  \AttributeTok{{-}{-}{-}{-}{-}{-}{-}}  \AttributeTok{{-}{-}{-}{-}{-}{-}{-}{-}{-}{-}{-}{-}{-}{-}{-}{-}{-}{-}{-}{-}{-}{-}{-}}
\ExtensionTok{health.health}\NormalTok{      GET      /health}
\ExtensionTok{notes.create\_note}\NormalTok{  POST     /notes/}
\ExtensionTok{notes.delete\_note}\NormalTok{  DELETE   /notes/}\OperatorTok{\textless{}}\NormalTok{int:note\_id}\OperatorTok{\textgreater{}}
\ExtensionTok{notes.get\_note}\NormalTok{     GET      /notes/}\OperatorTok{\textless{}}\NormalTok{int:note\_id}\OperatorTok{\textgreater{}}
\ExtensionTok{notes.list\_notes}\NormalTok{   GET      /notes/}
\ExtensionTok{static}\NormalTok{             GET      /static/}\OperatorTok{\textless{}}\NormalTok{path:filename}\OperatorTok{\textgreater{}}
\end{Highlighting}
\end{Shaded}

\end{codebox}

\noindent{}(\texttt{HEAD} and \texttt{OPTIONS} are hidden unless you add \texttt{-\/-all-methods}.)

\subsection*{Programmatic inspection}\phantomsection\label{32-routing:programmatic-inspection}

\begin{codebox}[short]

\begin{Shaded}
\begin{Highlighting}[]
\CommentTok{\# In a Flask shell or test:}
\ImportTok{from}\NormalTok{ app }\ImportTok{import}\NormalTok{ create\_app}
\NormalTok{app }\OperatorTok{=}\NormalTok{ create\_app()}

\CommentTok{\# Sorted by URL pattern for readability}
\NormalTok{rules }\OperatorTok{=} \BuiltInTok{sorted}\NormalTok{(app.url\_map.iter\_rules(), key}\OperatorTok{=}\KeywordTok{lambda}\NormalTok{ r: r.rule)}
\ControlFlowTok{for}\NormalTok{ rule }\KeywordTok{in}\NormalTok{ rules:}
\NormalTok{    methods }\OperatorTok{=} \StringTok{","}\NormalTok{.join(}\BuiltInTok{sorted}\NormalTok{(m }\ControlFlowTok{for}\NormalTok{ m }\KeywordTok{in}\NormalTok{ rule.methods }\ControlFlowTok{if}\NormalTok{ m }\KeywordTok{not} \KeywordTok{in}\NormalTok{ (}\StringTok{"HEAD"}\NormalTok{, }\StringTok{"OPTIONS"}\NormalTok{)))}
    \BuiltInTok{print}\NormalTok{(}\SpecialStringTok{f"  }\SpecialCharTok{\{}\NormalTok{methods}\SpecialCharTok{:10s\}}\SpecialStringTok{  }\SpecialCharTok{\{}\NormalTok{rule}\SpecialCharTok{.}\NormalTok{rule}\SpecialCharTok{:30s\}}\SpecialStringTok{  }\SpecialCharTok{\{}\NormalTok{rule}\SpecialCharTok{.}\NormalTok{endpoint}\SpecialCharTok{\}}\SpecialStringTok{"}\NormalTok{)}
\end{Highlighting}
\end{Shaded}

\end{codebox}

\subsection*{Building a route documentation endpoint}\phantomsection\label{32-routing:building-a-route-documentation-endpoint}

Some APIs expose their own route list for client developers:

\begin{codebox}[short]

\begin{Shaded}
\begin{Highlighting}[]
\CommentTok{\# In a Blueprint for administrative/meta endpoints:}
\ImportTok{from}\NormalTok{ flask }\ImportTok{import}\NormalTok{ current\_app}

\AttributeTok{@admin\_bp.route}\NormalTok{(}\StringTok{"/routes"}\NormalTok{)}
\KeywordTok{def}\NormalTok{ list\_routes():}
    \CommentTok{"""Return all registered routes as JSON. Development only."""}
    \ControlFlowTok{if} \KeywordTok{not}\NormalTok{ current\_app.debug:}
        \CommentTok{\# 404, not 403: in production, do not even confirm the endpoint exists}
        \ControlFlowTok{return}\NormalTok{ \{}\StringTok{"error"}\NormalTok{: }\StringTok{"Not found"}\NormalTok{\}, }\DecValTok{404}

\NormalTok{    routes }\OperatorTok{=}\NormalTok{ []}
    \ControlFlowTok{for}\NormalTok{ rule }\KeywordTok{in} \BuiltInTok{sorted}\NormalTok{(current\_app.url\_map.iter\_rules(), key}\OperatorTok{=}\KeywordTok{lambda}\NormalTok{ r: r.rule):}
\NormalTok{        routes.append(\{}
            \StringTok{"endpoint"}\NormalTok{: rule.endpoint,}
            \StringTok{"methods"}\NormalTok{: }\BuiltInTok{sorted}\NormalTok{(m }\ControlFlowTok{for}\NormalTok{ m }\KeywordTok{in}\NormalTok{ rule.methods }\ControlFlowTok{if}\NormalTok{ m }\KeywordTok{not} \KeywordTok{in}\NormalTok{ (}\StringTok{"HEAD"}\NormalTok{, }\StringTok{"OPTIONS"}\NormalTok{)),}
            \StringTok{"rule"}\NormalTok{: rule.rule,}
\NormalTok{        \})}
    \ControlFlowTok{return}\NormalTok{ \{}\StringTok{"routes"}\NormalTok{: routes\}}
\end{Highlighting}
\end{Shaded}

\end{codebox}

\setcounter{section}{8}
\section{How Flask Calls the Handler}\label{32-routing:329-how-flask-calls-the-handler}

Once routing identifies the endpoint and values, Flask calls the handler function. The full call sequence inside Flask is:

\begin{codebox}

\begin{Shaded}
\begin{Highlighting}[]
\CommentTok{\# Simplified from Flask\textquotesingle{}s full\_dispatch\_request method:}
\KeywordTok{def}\NormalTok{ full\_dispatch\_request(}\VariableTok{self}\NormalTok{):}
    \ControlFlowTok{try}\NormalTok{:}
        \CommentTok{\# 1. Run before\_request hooks}
\NormalTok{        rv }\OperatorTok{=} \VariableTok{self}\NormalTok{.preprocess\_request()}
        \ControlFlowTok{if}\NormalTok{ rv }\KeywordTok{is} \VariableTok{None}\NormalTok{:}
            \CommentTok{\# 2. Dispatch to the handler}
\NormalTok{            rv }\OperatorTok{=} \VariableTok{self}\NormalTok{.dispatch\_request()}
    \ControlFlowTok{except} \PreprocessorTok{Exception} \ImportTok{as}\NormalTok{ e:}
        \CommentTok{\# 3. Handle exceptions: the matching error handler produces rv}
\NormalTok{        rv }\OperatorTok{=} \VariableTok{self}\NormalTok{.handle\_user\_exception(e)}
    \CommentTok{\# 4. Make a response from the return value (and run after\_request hooks)}
    \ControlFlowTok{return} \VariableTok{self}\NormalTok{.finalize\_request(rv)}

\KeywordTok{def}\NormalTok{ dispatch\_request(}\VariableTok{self}\NormalTok{):}
    \CommentTok{\# Get the routing result from the request context}
\NormalTok{    rule }\OperatorTok{=}\NormalTok{ request.url\_rule       }\CommentTok{\# The matched Rule object}
\NormalTok{    view\_args }\OperatorTok{=}\NormalTok{ request.view\_args }\CommentTok{\# The extracted URL variables}

    \CommentTok{\# Look up the handler function}
    \ControlFlowTok{return} \VariableTok{self}\NormalTok{.ensure\_sync(}\VariableTok{self}\NormalTok{.view\_functions[rule.endpoint])(}\OperatorTok{**}\NormalTok{view\_args)}
\end{Highlighting}
\end{Shaded}

\end{codebox}

\noindent{}\texttt{view\_args} contains the extracted URL variables. For \texttt{GET\ /notes/42}, \texttt{view\_}\allowbreak{}\texttt{args\ =}\allowbreak{}\texttt{\ \{"note\_}\allowbreak{}\texttt{id":}\allowbreak{}\texttt{\ 42\}}, and Flask calls \texttt{get\_note(note\_id=}\allowbreak{}\texttt{42)}.

The \texttt{ensure\_sync()} call is Flask's support for async view functions (Python \texttt{async\ def}). For regular synchronous functions, it's a no-op passthrough.

\setcounter{section}{9}
\section{Routing in the Notes API}\label{32-routing:3210-routing-in-the-notes-api}

Let's examine the complete routing structure of the Notes API~--- as it will look once \hyperref[ch:33-logic-execution]{Chapter 33} adds the \texttt{PUT} route for updating a note:

\begin{codebox}

\begin{Shaded}
\begin{Highlighting}[]
\CommentTok{\# app/routes/notes.py}

\ImportTok{from}\NormalTok{ flask }\ImportTok{import}\NormalTok{ Blueprint}

\NormalTok{notes\_bp }\OperatorTok{=}\NormalTok{ Blueprint(}\StringTok{"notes"}\NormalTok{, }\VariableTok{\_\_name\_\_}\NormalTok{, url\_prefix}\OperatorTok{=}\StringTok{"/notes"}\NormalTok{)}

\AttributeTok{@notes\_bp.route}\NormalTok{(}\StringTok{"/"}\NormalTok{, methods}\OperatorTok{=}\NormalTok{[}\StringTok{"GET"}\NormalTok{])}
\KeywordTok{def}\NormalTok{ list\_notes():}
    \CommentTok{"""List all notes, newest first."""}
\NormalTok{    ...}

\AttributeTok{@notes\_bp.route}\NormalTok{(}\StringTok{"/"}\NormalTok{, methods}\OperatorTok{=}\NormalTok{[}\StringTok{"POST"}\NormalTok{])}
\KeywordTok{def}\NormalTok{ create\_note():}
    \CommentTok{"""Create a new note. Body: \{"content": "..."\}"""}
\NormalTok{    ...}

\AttributeTok{@notes\_bp.route}\NormalTok{(}\StringTok{"/\textless{}int:note\_id\textgreater{}"}\NormalTok{, methods}\OperatorTok{=}\NormalTok{[}\StringTok{"GET"}\NormalTok{])}
\KeywordTok{def}\NormalTok{ get\_note(note\_id: }\BuiltInTok{int}\NormalTok{):}
    \CommentTok{"""Get a specific note by ID."""}
\NormalTok{    ...}

\AttributeTok{@notes\_bp.route}\NormalTok{(}\StringTok{"/\textless{}int:note\_id\textgreater{}"}\NormalTok{, methods}\OperatorTok{=}\NormalTok{[}\StringTok{"PUT"}\NormalTok{])}
\KeywordTok{def}\NormalTok{ update\_note(note\_id: }\BuiltInTok{int}\NormalTok{):}
    \CommentTok{"""Update a note\textquotesingle{}s content. Body: \{"content": "..."\}"""}
\NormalTok{    ...}

\AttributeTok{@notes\_bp.route}\NormalTok{(}\StringTok{"/\textless{}int:note\_id\textgreater{}"}\NormalTok{, methods}\OperatorTok{=}\NormalTok{[}\StringTok{"DELETE"}\NormalTok{])}
\KeywordTok{def}\NormalTok{ delete\_note(note\_id: }\BuiltInTok{int}\NormalTok{):}
    \CommentTok{"""Delete a note."""}
\NormalTok{    ...}
\end{Highlighting}
\end{Shaded}

\end{codebox}

\begin{codebox}[short]

\begin{Shaded}
\begin{Highlighting}[]
\CommentTok{\# app/\_\_init\_\_.py}
\KeywordTok{def}\NormalTok{ create\_app():}
\NormalTok{    app }\OperatorTok{=}\NormalTok{ Flask(}\VariableTok{\_\_name\_\_}\NormalTok{)}
    \CommentTok{\# ...}
    \ImportTok{from}\NormalTok{ app.routes.notes }\ImportTok{import}\NormalTok{ notes\_bp}
\NormalTok{    app.register\_blueprint(notes\_bp)}
    \ControlFlowTok{return}\NormalTok{ app}
\end{Highlighting}
\end{Shaded}

\end{codebox}

\noindent{}The resulting URL map:

\begin{outputbox}[short]
\begin{Verbatim}[fontsize=\codesize, breaklines, breakanywhere, breaksymbolleft={\tiny\textcolor{muted}{$\hookrightarrow$}}]
Method    URL                    Handler
GET       /notes/                notes.list_notes
POST      /notes/                notes.create_note
GET       /notes/<int:note_id>   notes.get_note
PUT       /notes/<int:note_id>   notes.update_note
DELETE    /notes/<int:note_id>   notes.delete_note
\end{Verbatim}
\end{outputbox}

\noindent{}This design is \textbf{resource-oriented}: URLs identify resources (\texttt{/notes/}, \texttt{/notes/42}), and HTTP methods identify actions (\texttt{GET} = read, \texttt{POST} = create, \texttt{PUT} = replace, \texttt{DELETE} = remove). A note has only one field a client can set, so \texttt{PUT} (replace the note) and \texttt{PATCH} (change some fields) would do the same thing here; with more fields, the difference matters. This is RESTful API design, and it maps cleanly onto Flask's routing primitives.

\phantomsection\label{32-routing:key-takeaways}
\begin{takeaways}

\begin{itemize}
\tightlist
\item
  \textbf{\texttt{@app.route()}} is a decorator that calls \texttt{add\_url\_rule()}, which creates a Werkzeug \texttt{Rule} and adds it to the app's \texttt{url\_map}. Registration happens at import time.
\item
  \textbf{The \texttt{url\_map}} is a collection of \texttt{Rule} objects. Each \texttt{Rule} contains the URL pattern, the allowed HTTP methods, and the endpoint name; Werkzeug compiles them all into one state machine for matching.
\item
  \textbf{URL matching} finds the most specific rule that fits the path. If no Rule matches the path, 404 is raised. If a Rule matches the path but not the method, 405 is raised.
\item
  \textbf{Converters} (\texttt{\textless{}int:note\_id\textgreater{}}) both constrain what the URL segment must look like and convert the matched string to the appropriate Python type. \texttt{/notes/abc} returns 404 when the route uses \texttt{\textless{}int:note\_id\textgreater{}}.
\item
  \textbf{Trailing slash rules}: routes with a trailing slash redirect requests without one (308); routes without a trailing slash return 404 for requests with one. API clients should use the exact URLs.
\item
  \textbf{Blueprints} collect routes separately and are merged into the app's \texttt{url\_map} at \texttt{register\_blueprint()} time. Endpoint names are prefixed with the Blueprint name.
\item
  \textbf{\texttt{url\_for()}} builds URLs from endpoint names. Always use \texttt{url\_for()} instead of hardcoding paths so that URL changes propagate automatically.
\end{itemize}

\end{takeaways}

\unnumberedsection{false}{Exercises}\label{32-routing:exercises}

\begin{enumerate}
\def\labelenumi{\arabic{enumi}.}
\tightlist
\item
  \textbf{Predict.} Request \texttt{/notes} (no trailing slash) with \texttt{curl\ -i}. What status and header come back, and why is that a problem for a \texttt{POST}?
\item
  \textbf{Break and fix.} Add \texttt{@notes\_}\allowbreak{}\texttt{bp.}\allowbreak{}\texttt{route("/}\allowbreak{}\texttt{\textless{}note\_}\allowbreak{}\texttt{id\textgreater{}")} (no converter) next to the \texttt{int} route. Which one handles \texttt{/notes/5}? Remove the ambiguity.
\item
  \textbf{Extend.} Print the application's URL map (\texttt{flask\ -}\allowbreak{}\texttt{-}\allowbreak{}\texttt{app\ \textquotesingle{}app:}\allowbreak{}\texttt{create\_}\allowbreak{}\texttt{app()\textquotesingle{}\ routes}) and compare it with the endpoints listed in section 32.10. Is there a rule you did not expect?
\item
  \textbf{Explain.} Why should code build URLs with \texttt{url\_for()} instead of writing paths as strings?
\end{enumerate}

\noindent{}Solutions and hints are in the companion repository, in \texttt{exercises/}\allowbreak{}\texttt{chapter-32.md}. Try each exercise before reading its solution: the point is the thinking, not the answer.

\unnumberedsection{false}{What's Next}\label{32-routing:whats-next}

The routing layer has identified the handler function and passed in the extracted URL variables. The handler now executes its logic: fetching data from the database, applying business rules, computing results. That execution is the subject of the next chapter.

\noindent{}→ \hyperref[ch:33-logic-execution]{Chapter 33}~--- Logic Execution
